\section{Expressions and processes}
\label{sec:expr-proc}

% Figures in this section: syntax of expressions (\cref{fig:syntax-expressions}),
% run-time syntax (\cref{fig:syntax-runtime}), reduction (expressions:
% \cref{fig:expression-reduction}, processes: \cref{fig:process-reduction}),
% process congruence (\cref{fig:configuration-congruence}).

\begin{figure}
  \begin{align*}
    \Dir \grmeq
      & \DirLeft \grmor \DirRight
        % \grmor \DirNone
        \tag{Directions} \\
    \C \grmeq
      & \CUnit
        \grmor \CNew
        \grmor \CFork
        \grmor \CTerm 
        \grmor \CWait
        \grmor \CSend
        \grmor \CRecv
    \\\grmor&
     \CSplit \grmor \CSSplit
    \grmor \CDrop \grmor \CAcq
        \grmor \CSelect k
        \grmor \CBranch
        \tag{Constants} \\
    \E \grmeq
      & \C \grmor \EVar \grmor \EAbs\EVar\E \grmor \EAppD{\E}{\E}{\Dir} 
        \grmor \ESeq\E\E
        \grmor \EPair\E\E
        \grmor \ELetPair{\EVar}{\EVar}{\E}{\E} \\
    \grmor& \EInj i \E
        %\grmor \ECase\E{\EVar}{\E}{\EVar}{\E}
        \grmor \ECase \E {\ECaseBranch* i \EVar \E}
        \grmor \EMu[\EVar] \E
      % & \grmor \EMatch{\E}{\ell:\E_\ell}{\ell\in L}
      % \grmor \EFork\E \grmor \ELet\EVar\E\E
        \tag{Expressions}
  \end{align*}
  \caption{Syntax of expressions}
  \label{fig:syntax-expressions}
\end{figure}

Figure \cref{fig:syntax-expressions} contains the syntax of constants and expressions. 
Constants $\C$ include the unit value and functional constants for channel creation, thread spawning, 
and manipulating with existing channels. They are better explained together with their type later.  

Expressions $\E$ contain constants, variables, lambda abstraction and application, sequencing, 
pair introduction and pair elimination, sum injection and case analysis, and recursive terms. 

We distinguish between left ($\DirLeft$) and right ($\DirRight$) application. 
The direction determines the order of evaluation: either the left or right term is evaluated first.
Throughout this paper, an omitted direction annotation is understood to denote left-directed application.

\begin{figure}[t!]
  \begin{align*}
    \U,\V \grmeq
    & \C
      \grmor \EVar
      \grmor \EAbs \EVar\E
      \grmor \EInj{i} \V
      \grmor \EPair \U\V
      % \grmor \CSend \V 
              \tag{Values}
    \\
    \EC \grmeq&
              \ECEmpty
              \grmor \EAppD \EC\E\DirLeft
              \grmor \EAppD \V\EC\DirLeft
              \grmor \EAppD \E\EC\DirRight
              \grmor \EAppD \EC\V\DirRight
              \grmor \ESeq \EC\E
              \grmor \EPair \EC\E
              \grmor \EPair \V\EC
              \\ \grmor
    & \EInj{i}\EC \grmor \ELet {\EPair \EVar\EVar}\EC\E
              \grmor \ECase \EC {\ECaseBranch* i \EVar \E}
              \tag{Evaluation contexts}
    \\
    \BindGroup \grmeq
    & (\EVar \grmor \BSep)^*
      \tag{Binder groups}
    \\
    \Bind \grmeq
    & \EVar^*
      \tag{Variable sequences}
    \\
    \Proc \grmeq
    & \PExp\E
      \grmor \PPar\Proc\Proc
      \grmor \PNew{\BindGroup}{\BindGroup}{\Proc}
      \tag{Processes}
    \\
    \FC \grmeq
    & \PExp \EC
      \tag{Thread evaluation contexts}
    \\
    \CC \grmeq
    & \ECEmpty
      \grmor \PPar{\CC}{\Proc}
      \grmor \PNew \BindGroup\BindGroup \CC
      \tag{Process contexts}
  \end{align*}
  \caption{Run-time syntax}
  \label{fig:syntax-runtime}
\end{figure}

Figure \cref{fig:syntax-runtime} contains the run time syntax that evaluates expressions inside processes. 
The values $\U,\V$ constitute a subset of expressions that are irreducible. Evaluation contexts $\EC$ define where expression reduction may take the next step. 

Binder groups $\BindGroup$ are flat lists of channel names $\EVar$ and
separators $\BSep$. The metavariable $\Bind$ ranges over separator-free
variable sequences. We write $\BEmp$ for the empty list and
$\BWithSep{\Bind}{\BindGroup}$ for the concatenation of $\Bind$, $\BSep$, and
$\BindGroup$. The paper only creates groups where two separators never occur
side-by-side without an intervening channel name, but the syntax does not
enforce this invariant. Processes $\Proc$ include a single expression, parallel
composition of processes, and channel introduction with two binder groups.
Finally, we lift the evaluation contexts to processes by defining thread
contexts $\FC$ and process contexts $\CC$.

Binder sequence concatenation is written by juxtaposition with unit $\BEmp$.

% EXP EVALUATION

\begin{figure}[t!]
  \begin{mathpar}
    \RuleExpRedApp \and
    \RuleExpRedSeq \and
    \RuleExpRedPairElim \and
    \RuleExpRedUnfold \and
    \RuleExpRedVariantElim \and
    \RuleExpRedCtx
  \end{mathpar}
  \caption{Expression reduction ($\E \EReduce \E$)}
  \label{fig:expression-reduction}
\end{figure}

Processes $\Proc$ inlcude a single expression, parallel composition of the former, and channel introduction with two binder groups. 
Finally, we can lift the evaluation contexts to processes by defining thread contexts $\FC$ and process contexts $\CC$. 

\begin{figure}[t!]
  \begin{mathpar}
    \RuleProcRedExp \and
    \RuleProcRedNew \and
    \RuleProcRedFork \and
    \RuleProcRedCom \and
    \RuleProcRedLSplit \and
    \RuleProcRedRSplit \and
    \RuleProcRedDiscard \and
    \RuleProcRedDrop \and
    \RuleProcRedAcquire \and
    \RuleProcRedClose \and
    \RuleProcRedChoice \and
    \RuleProcRedPar \and
    \RuleProcRedBind \and
    \RuleProcRedStruct
  \end{mathpar}
  \caption{Process reduction ($\Proc \PReduce \Proc$)}
  \label{fig:process-reduction}
\end{figure}


%%% Local Variables:
%%% mode: latex
%%% TeX-master: "../popl2027.tex"
%%% End:
